% Constellation diagrams for the modulation section of the radio
% chapter.
%
% Every one of these is the same picture with different dots on it: a
% real and an imaginary axis, tick marks at +/-1, and a symbol at each
% point the modulator can produce. Keeping the frame in one macro means
% the four diagrams are actually comparable, which is the whole reason
% the chapter shows them one after another.
%
% \constframe{<half-width>}  draws the axes and their labels.
% \constsymbol{re}{im}       marks one symbol.
% \constcircle{r}            the unit circle, for the phase-only schemes.

\newcommand{\constframe}[1]{
  \draw[->] (-#1-0.35,0) -- (#1+0.45,0) node[anchor=north east] {$\Re$};
  \draw[->] (0,-#1-0.35) -- (0,#1+0.45) node[anchor=north east] {$\Im$};
  \foreach \t in {-1,1} {
    \draw (\t,-0.08) -- (\t,0.08);
    \draw (-0.08,\t) -- (0.08,\t);
  }
}

% A symbol. Red, because in this book red is the thing the figure is
% about, and here the symbols are the whole content.
% The braces let the arguments be expressions such as cos(45), which
% TikZ otherwise tries to read as a node name because of the brackets.
\newcommand{\constsymbol}[2]{
  \fill[red] ({#1},{#2}) circle (0.075);
}

\newcommand{\constcircle}[1]{
  \draw[gray, dotted, thick] (0,0) circle (#1);
}
